Dense metric depth estimation from a single monocular image is fundamental for autonomous robotics, micro-aerial vehicles (MAVs), and spatial computing. However, current state-of-the-art metric foundation models rely on massive transformer backbones ($>300$\,M parameters, $>1$\,GB memory), precluding real-time deployment on compute-constrained edge platforms. Conversely, lightweight models trained from scratch frequently collapse into blurry distance representations and suffer from severe focal-depth ambiguity. We present \textbf{Dioptra-DINO}, a compact, geometry-coupled architecture ($27.51$\,M parameters, $105.05$\,MB weights) designed to operate within an edge-robotics compute envelope. Dioptra-DINO couples a self-supervised DINOv2-Small vision backbone with analytical projective camera geometry: (i)~\textit{Trivision Ray Positional Encoding}, modulating transformer token representations via affine Feature-wise Linear Modulation (FiLM) driven by continuous optical ray triplets unprojected via Cramer's rule closed-form inversion; (ii)~\textit{Angular Residual Attention (ARA)}, a geometric attention bias that penalizes divergent optical rays to sharpen object silhouettes and structural depth boundaries; (iii)~\textit{Decoupled Metric Scale \& Shape Supervision}, combining raw physical metric SiLog, explicit log-median scale consistency, and Multi-Scale 3D Virtual Normal Loss ($\mathcal{L}_{\text{vnl}}$) to enforce planar architectural surface orientation; and (iv)~\textit{Dynamic Optical Pinhole Augmentation}, enforcing camera-intrinsic equivariance through continuous focal sampling ($f \in [0.75f_0, 1.5f_0]$). Evaluated on unseen cross-environment transfer across a continuous 200-image trajectory in the held-out \textit{abandonedfactory/Easy/P010} benchmark, Dioptra-DINO achieves \textbf{$0.0555$ AbsRel} ($5.55\%$ relative error), \textbf{$97.06\%$ $\delta_1 < 1.25$ threshold accuracy}, \textbf{$2.396$\,m RMSE}, and an absolute median metric scale ratio of \textbf{$1.0003$} (only $0.03\%$ scale error) without any test-time scale alignment, outperforming the leading lightweight foundation baseline Depth Anything V2-Small ($0.0964$ AbsRel, $90.78\%$ $\delta_1$ under oracle per-frame affine alignment using ground truth) by $42.4\%$ and the camera-conditioned foundation model Metric3D ViT-Small ($0.3269$ AbsRel, $26.52\%$ $\delta_1$) by $83.0\%$ while operating up to $21.1\times$ faster ($38.0$\,FPS vs.\ $1.8$\,FPS). Synthetic camera field-of-view sweeps ($50^\circ \text{--} 100^\circ$) across multi-frame sequences ($N=20$) demonstrate projective equivariance, maintaining scale calibration within $[-1.4\%, +3.0\%]$ ($0.9863\text{--}1.0295$) across the wide-angle operational range ($73.7^\circ \text{--} 100^\circ$) with $>96.4\%$ $\delta_1$ accuracy. Conversely, zero-shot evaluation on the real-world indoor benchmark NYU Depth V2 reveals an expected metric scale expansion ($1.28\times$ scale ratio; AbsRel $0.4867$, $\delta_1 = 46.25\%$), establishing that while projective ray unprojection guarantees optical equivariance across changing focal lengths, open-world absolute metric scale generalization strictly requires diverse multi-domain pretraining. Operating at steady-state forward throughput of \textbf{$38.0 \pm 2.2$\,FPS} ($26.3 \pm 1.5$\,ms forward latency; $28.3 \pm 1.4$\,FPS / $35.3 \pm 1.8$\,ms full pipeline across 100 repeated trials) on an Apple Silicon M3 GPU under unquantized FP32 with peak working memory under $210$\,MB, Dioptra-DINO demonstrates that explicit geometric inductive biases enable a compact model to deliver robust metric depth within an edge-compatible computational envelope. Code, raw per-frame evaluation logs, and model checkpoints are available at \url{https://github.com/SeranomTheGreat/dioptra}.

